\documentclass[]{article}
\usepackage{lmodern}
\usepackage{amssymb,amsmath}
\usepackage{ifxetex,ifluatex}
\usepackage{fixltx2e} % provides \textsubscript
\ifnum 0\ifxetex 1\fi\ifluatex 1\fi=0 % if pdftex
  \usepackage[T1]{fontenc}
  \usepackage[utf8]{inputenc}
\else % if luatex or xelatex
  \ifxetex
    \usepackage{mathspec}
  \else
    \usepackage{fontspec}
  \fi
  \defaultfontfeatures{Ligatures=TeX,Scale=MatchLowercase}
\fi
% use upquote if available, for straight quotes in verbatim environments
\IfFileExists{upquote.sty}{\usepackage{upquote}}{}
% use microtype if available
\IfFileExists{microtype.sty}{%
\usepackage{microtype}
\UseMicrotypeSet[protrusion]{basicmath} % disable protrusion for tt fonts
}{}
\usepackage[margin=1in]{geometry}
\usepackage{hyperref}
\hypersetup{unicode=true,
            pdftitle={S3: The direct measurement of an optimal offset},
            pdfborder={0 0 0},
            breaklinks=true}
\urlstyle{same}  % don't use monospace font for urls
\usepackage{color}
\usepackage{fancyvrb}
\newcommand{\VerbBar}{|}
\newcommand{\VERB}{\Verb[commandchars=\\\{\}]}
\DefineVerbatimEnvironment{Highlighting}{Verbatim}{commandchars=\\\{\}}
% Add ',fontsize=\small' for more characters per line
\usepackage{framed}
\definecolor{shadecolor}{RGB}{248,248,248}
\newenvironment{Shaded}{\begin{snugshade}}{\end{snugshade}}
\newcommand{\AlertTok}[1]{\textcolor[rgb]{0.94,0.16,0.16}{#1}}
\newcommand{\AnnotationTok}[1]{\textcolor[rgb]{0.56,0.35,0.01}{\textbf{\textit{#1}}}}
\newcommand{\AttributeTok}[1]{\textcolor[rgb]{0.77,0.63,0.00}{#1}}
\newcommand{\BaseNTok}[1]{\textcolor[rgb]{0.00,0.00,0.81}{#1}}
\newcommand{\BuiltInTok}[1]{#1}
\newcommand{\CharTok}[1]{\textcolor[rgb]{0.31,0.60,0.02}{#1}}
\newcommand{\CommentTok}[1]{\textcolor[rgb]{0.56,0.35,0.01}{\textit{#1}}}
\newcommand{\CommentVarTok}[1]{\textcolor[rgb]{0.56,0.35,0.01}{\textbf{\textit{#1}}}}
\newcommand{\ConstantTok}[1]{\textcolor[rgb]{0.00,0.00,0.00}{#1}}
\newcommand{\ControlFlowTok}[1]{\textcolor[rgb]{0.13,0.29,0.53}{\textbf{#1}}}
\newcommand{\DataTypeTok}[1]{\textcolor[rgb]{0.13,0.29,0.53}{#1}}
\newcommand{\DecValTok}[1]{\textcolor[rgb]{0.00,0.00,0.81}{#1}}
\newcommand{\DocumentationTok}[1]{\textcolor[rgb]{0.56,0.35,0.01}{\textbf{\textit{#1}}}}
\newcommand{\ErrorTok}[1]{\textcolor[rgb]{0.64,0.00,0.00}{\textbf{#1}}}
\newcommand{\ExtensionTok}[1]{#1}
\newcommand{\FloatTok}[1]{\textcolor[rgb]{0.00,0.00,0.81}{#1}}
\newcommand{\FunctionTok}[1]{\textcolor[rgb]{0.00,0.00,0.00}{#1}}
\newcommand{\ImportTok}[1]{#1}
\newcommand{\InformationTok}[1]{\textcolor[rgb]{0.56,0.35,0.01}{\textbf{\textit{#1}}}}
\newcommand{\KeywordTok}[1]{\textcolor[rgb]{0.13,0.29,0.53}{\textbf{#1}}}
\newcommand{\NormalTok}[1]{#1}
\newcommand{\OperatorTok}[1]{\textcolor[rgb]{0.81,0.36,0.00}{\textbf{#1}}}
\newcommand{\OtherTok}[1]{\textcolor[rgb]{0.56,0.35,0.01}{#1}}
\newcommand{\PreprocessorTok}[1]{\textcolor[rgb]{0.56,0.35,0.01}{\textit{#1}}}
\newcommand{\RegionMarkerTok}[1]{#1}
\newcommand{\SpecialCharTok}[1]{\textcolor[rgb]{0.00,0.00,0.00}{#1}}
\newcommand{\SpecialStringTok}[1]{\textcolor[rgb]{0.31,0.60,0.02}{#1}}
\newcommand{\StringTok}[1]{\textcolor[rgb]{0.31,0.60,0.02}{#1}}
\newcommand{\VariableTok}[1]{\textcolor[rgb]{0.00,0.00,0.00}{#1}}
\newcommand{\VerbatimStringTok}[1]{\textcolor[rgb]{0.31,0.60,0.02}{#1}}
\newcommand{\WarningTok}[1]{\textcolor[rgb]{0.56,0.35,0.01}{\textbf{\textit{#1}}}}
\usepackage{graphicx,grffile}
\makeatletter
\def\maxwidth{\ifdim\Gin@nat@width>\linewidth\linewidth\else\Gin@nat@width\fi}
\def\maxheight{\ifdim\Gin@nat@height>\textheight\textheight\else\Gin@nat@height\fi}
\makeatother
% Scale images if necessary, so that they will not overflow the page
% margins by default, and it is still possible to overwrite the defaults
% using explicit options in \includegraphics[width, height, ...]{}
\setkeys{Gin}{width=\maxwidth,height=\maxheight,keepaspectratio}
\IfFileExists{parskip.sty}{%
\usepackage{parskip}
}{% else
\setlength{\parindent}{0pt}
\setlength{\parskip}{6pt plus 2pt minus 1pt}
}
\setlength{\emergencystretch}{3em}  % prevent overfull lines
\providecommand{\tightlist}{%
  \setlength{\itemsep}{0pt}\setlength{\parskip}{0pt}}
\setcounter{secnumdepth}{0}
% Redefines (sub)paragraphs to behave more like sections
\ifx\paragraph\undefined\else
\let\oldparagraph\paragraph
\renewcommand{\paragraph}[1]{\oldparagraph{#1}\mbox{}}
\fi
\ifx\subparagraph\undefined\else
\let\oldsubparagraph\subparagraph
\renewcommand{\subparagraph}[1]{\oldsubparagraph{#1}\mbox{}}
\fi

%%% Use protect on footnotes to avoid problems with footnotes in titles
\let\rmarkdownfootnote\footnote%
\def\footnote{\protect\rmarkdownfootnote}

%%% Change title format to be more compact
\usepackage{titling}

% Create subtitle command for use in maketitle
\newcommand{\subtitle}[1]{
  \posttitle{
    \begin{center}\large#1\end{center}
    }
}

\setlength{\droptitle}{-2em}

  \title{S3: The direct measurement of an optimal offset}
    \pretitle{\vspace{\droptitle}\centering\huge}
  \posttitle{\par}
    \author{true \\ true}
    \preauthor{\centering\large\emph}
  \postauthor{\par}
      \predate{\centering\large\emph}
  \postdate{\par}
    \date{2019-03-04}

\usepackage{booktabs}
\usepackage{longtable}
\usepackage{array}
\usepackage{multirow}
\usepackage{wrapfig}
\usepackage{float}
\usepackage{colortbl}
\usepackage{pdflscape}
\usepackage{tabu}
\usepackage{threeparttable}
\usepackage{threeparttablex}
\usepackage[normalem]{ulem}
\usepackage{makecell}
\usepackage{xcolor}

\begin{document}
\maketitle

{
\setcounter{tocdepth}{3}
\tableofcontents
}
Consider the Mamre interaction data set (Pauw, personal communications)
replicated in table A1. This data summarizes the fitness effects of
pollination-based interactions between individuals of the long-mouthed
fly species \emph{M. Longirostris} and the long-tubed flower \emph{L.
anceps}.

\begin{verbatim}
## Loading required package: knitr
\end{verbatim}

\begin{verbatim}
## Loading required package: kableExtra
\end{verbatim}

\begin{verbatim}
## Loading required package: xlsx
\end{verbatim}

\begin{table}[t]

\caption{\label{tab:unnamed-chunk-1}Measurements of fitness and morphological differences}
\centering
\begin{tabular}{r|r|r}
\hline
nectar consumed & pollen grains deposited & tongue length - tube length\\
\hline
0.480 & 10 & -6.5\\
\hline
1.702 & 8 & -7.0\\
\hline
1.344 & 7 & -8.5\\
\hline
0.528 & 6 & -13.0\\
\hline
0.000 & 2 & -17.0\\
\hline
0.066 & 6 & -10.5\\
\hline
0.000 & 44 & -11.0\\
\hline
0.414 & 216 & -6.0\\
\hline
0.480 & 3 & 3.0\\
\hline
2.808 & 71 & -3.0\\
\hline
0.378 & 106 & -3.0\\
\hline
0.288 & 127 & -4.5\\
\hline
0.000 & 0 & -6.0\\
\hline
1.656 & 129 & -6.5\\
\hline
0.000 & 67 & -14.5\\
\hline
0.190 & 61 & 0.5\\
\hline
0.174 & 8 & -1.0\\
\hline
0.330 & 117 & -6.5\\
\hline
0.000 & 84 & -7.5\\
\hline
0.000 & 79 & -12.0\\
\hline
0.000 & 205 & -17.0\\
\hline
0.000 & 33 & -18.5\\
\hline
0.414 & 2 & -2.0\\
\hline
1.470 & 104 & -10.0\\
\hline
0.414 & 0 & 6.5\\
\hline
1.950 & 26 & 4.5\\
\hline
0.294 & 1 & 1.0\\
\hline
1.950 & 0 & -0.5\\
\hline
1.584 & 42 & -1.0\\
\hline
0.960 & 108 & -10.0\\
\hline
3.328 & 115 & -12.0\\
\hline
0.000 & 91 & -15.0\\
\hline
0.426 & 122 & -16.5\\
\hline
0.870 & 9 & 3.5\\
\hline
1.470 & 12 & 0.5\\
\hline
1.360 & 4 & -1.5\\
\hline
0.702 & 61 & -2.0\\
\hline
1.224 & 24 & -3.5\\
\hline
0.486 & 83 & -4.0\\
\hline
1.430 & 62 & -11.0\\
\hline
1.152 & 298 & -8.5\\
\hline
0.240 & 209 & -4.0\\
\hline
1.344 & 21 & -5.0\\
\hline
0.754 & 1 & 5.5\\
\hline
0.550 & 5 & 5.0\\
\hline
0.702 & 6 & 4.5\\
\hline
1.360 & 1 & 4.0\\
\hline
0.414 & 5 & 3.5\\
\hline
0.784 & 28 & -1.5\\
\hline
0.000 & 87 & -11.0\\
\hline
0.000 & 63 & -17.5\\
\hline
1.064 & 0 & 5.0\\
\hline
1.520 & 4 & 1.0\\
\hline
0.000 & 93 & -7.5\\
\hline
0.550 & 127 & -4.0\\
\hline
0.990 & 56 & -8.5\\
\hline
0.000 & 32 & -4.0\\
\hline
2.494 & 16 & 2.0\\
\hline
2.242 & 23 & -7.0\\
\hline
0.414 & 0 & 10.0\\
\hline
0.352 & 3 & 9.5\\
\hline
0.550 & 0 & 8.0\\
\hline
0.870 & 32 & 5.0\\
\hline
0.480 & 62 & 2.5\\
\hline
3.430 & 2 & 4.0\\
\hline
1.254 & 36 & -2.0\\
\hline
2.214 & 1 & 5.5\\
\hline
1.008 & 57 & -3.5\\
\hline
1.054 & 19 & 3.0\\
\hline
0.638 & 12 & -2.0\\
\hline
0.294 & 24 & 13.0\\
\hline
1.290 & 168 & -4.5\\
\hline
3.774 & 105 & 8.0\\
\hline
3.430 & 16 & 5.5\\
\hline
1.650 & 129 & -7.0\\
\hline
1.584 & 42 & 18.5\\
\hline
0.960 & 12 & 17.0\\
\hline
0.870 & 2 & 17.0\\
\hline
1.824 & 1 & 18.0\\
\hline
2.214 & 48 & 17.5\\
\hline
0.870 & 39 & 13.0\\
\hline
\end{tabular}
\end{table}

To measure the optimal offset from these data, we simply minimize the
sum of the square errors between the offset-matching fitness surfaces
and the fitness of each species, constraining the \(\delta\)'s for each
surface to be the same. That is, denoting \(\hat W_i^j\) and \(z_i^j\)
the fitness and trait value of species \(i\) during observation \(j\),
the data can be written as

\(D=\{\hat W^i,\Delta_i\}_{i=1}^n=\{\{\hat W_1^i,\hat W_2^i\},\{z_1^i-z_2^i\}\}_{i=1}^n\)

And hence we minimize the quantity

\[L^2(W_1,W_2|D)\equiv\sum_{i=1}^n\left(\hat W_1^i-W_1(\Delta_i)\right)^2+\left(\hat W_2^i-W_2(\Delta_i)\right)^2.\]

The fitness functions can be written as

\[W_i(z_j-z_i|c_i,B_i,\delta)=c_i\exp\left(-\frac{B_i}{2}(\delta+z_j-z_i)^2\right)\]

and hence the goal is to compute

\[{\arg\min}_{c_1,c_2,B_1,B_2,\delta}(L^2).\]

This can be performed numerically. For the data in table A1, we find

\begin{Shaded}
\begin{Highlighting}[]
\NormalTok{W <-}\StringTok{ }\ControlFlowTok{function}\NormalTok{(D,c,B,d)\{}
\NormalTok{  f <-}\StringTok{ }\NormalTok{c}\OperatorTok{*}\KeywordTok{exp}\NormalTok{(}\OperatorTok{-}\NormalTok{B}\OperatorTok{*}\NormalTok{(d}\OperatorTok{+}\NormalTok{D)}\OperatorTok{/}\DecValTok{2}\NormalTok{)}
  \KeywordTok{return}\NormalTok{(f)}
\NormalTok{\}}

\CommentTok{# collect data}
\NormalTok{W1 <-}\StringTok{ }\NormalTok{mamre}\OperatorTok{$}\StringTok{`}\DataTypeTok{nectar consumed}\StringTok{`}
\NormalTok{W2 <-}\StringTok{ }\NormalTok{mamre}\OperatorTok{$}\StringTok{`}\DataTypeTok{pollen grains deposited}\StringTok{`}
\NormalTok{D  <-}\StringTok{ }\NormalTok{mamre}\OperatorTok{$}\StringTok{`}\DataTypeTok{tongue length - tube length}\StringTok{`}
\NormalTok{n <-}\StringTok{ }\KeywordTok{dim}\NormalTok{(mamre)[}\DecValTok{1}\NormalTok{]}
  
\NormalTok{L <-}\StringTok{ }\ControlFlowTok{function}\NormalTok{(X)\{}
  
  \CommentTok{# unpack parameters}
\NormalTok{  c1 <-}\StringTok{ }\KeywordTok{exp}\NormalTok{(X[}\DecValTok{1}\NormalTok{])}
\NormalTok{  c2 <-}\StringTok{ }\KeywordTok{exp}\NormalTok{(X[}\DecValTok{2}\NormalTok{])}
\NormalTok{  B1 <-}\StringTok{ }\KeywordTok{exp}\NormalTok{(X[}\DecValTok{3}\NormalTok{])}
\NormalTok{  B2 <-}\StringTok{ }\KeywordTok{exp}\NormalTok{(X[}\DecValTok{4}\NormalTok{])}
\NormalTok{  d  <-}\StringTok{ }\KeywordTok{exp}\NormalTok{(X[}\DecValTok{5}\NormalTok{])}
  
\NormalTok{  l <-}\StringTok{ }\DecValTok{0}
  \CommentTok{# accumulate errors across M. longir}
  \ControlFlowTok{for}\NormalTok{(i }\ControlFlowTok{in} \DecValTok{1}\OperatorTok{:}\NormalTok{n)\{}
\NormalTok{    l <-}\StringTok{ }\NormalTok{l }\OperatorTok{+}\StringTok{ }\NormalTok{( }\KeywordTok{W}\NormalTok{(}\OperatorTok{-}\NormalTok{D[i],c1,B1,d) }\OperatorTok{-}\StringTok{ }\NormalTok{W1[i] )}\OperatorTok{^}\DecValTok{2}
\NormalTok{  \}}
  
  \CommentTok{# accumulate errors across L. anceps}
  \ControlFlowTok{for}\NormalTok{(i }\ControlFlowTok{in} \DecValTok{1}\OperatorTok{:}\NormalTok{n)\{}
\NormalTok{    l <-}\StringTok{ }\NormalTok{l }\OperatorTok{+}\StringTok{ }\NormalTok{( }\KeywordTok{W}\NormalTok{(D[i],c2,B2,d) }\OperatorTok{-}\StringTok{ }\NormalTok{W2[i] )}\OperatorTok{^}\DecValTok{2}
\NormalTok{  \}}
  
  \KeywordTok{return}\NormalTok{(}\KeywordTok{sqrt}\NormalTok{(l))}
  
\NormalTok{\}}

\NormalTok{mins <-}\StringTok{ }\OtherTok{NULL}
\NormalTok{vals <-}\StringTok{ }\OtherTok{NULL}
\ControlFlowTok{for}\NormalTok{(i }\ControlFlowTok{in} \DecValTok{1}\OperatorTok{:}\DecValTok{1000}\NormalTok{)\{}
\NormalTok{  out <-}\StringTok{ }\KeywordTok{optim}\NormalTok{(}\KeywordTok{rnorm}\NormalTok{(}\DecValTok{5}\NormalTok{),L)}
\NormalTok{  mins <-}\StringTok{ }\KeywordTok{rbind}\NormalTok{(mins,out}\OperatorTok{$}\NormalTok{par)}
\NormalTok{  vals <-}\StringTok{ }\KeywordTok{c}\NormalTok{(vals,out}\OperatorTok{$}\NormalTok{value)}
\NormalTok{\}}
\NormalTok{i <-}\StringTok{ }\KeywordTok{which.min}\NormalTok{(vals)}
\NormalTok{argmin <-}\StringTok{ }\KeywordTok{exp}\NormalTok{(mins[i,])}

\NormalTok{c1 <-}\StringTok{ }\NormalTok{argmin[}\DecValTok{1}\NormalTok{]}
\NormalTok{c2 <-}\StringTok{ }\NormalTok{argmin[}\DecValTok{2}\NormalTok{]}
\NormalTok{B1 <-}\StringTok{ }\NormalTok{argmin[}\DecValTok{3}\NormalTok{]}
\NormalTok{B2 <-}\StringTok{ }\NormalTok{argmin[}\DecValTok{4}\NormalTok{]}
\NormalTok{d  <-}\StringTok{ }\NormalTok{argmin[}\DecValTok{5}\NormalTok{]}

\NormalTok{fits <-}\StringTok{ }\KeywordTok{data.frame}\NormalTok{(}\KeywordTok{c}\NormalTok{(}\StringTok{"M. longirostris"}\NormalTok{,}\StringTok{"L. anceps"}\NormalTok{),}\KeywordTok{c}\NormalTok{(c1,c2),}\KeywordTok{c}\NormalTok{(B1,B2),}\KeywordTok{c}\NormalTok{(d,d))}
\KeywordTok{colnames}\NormalTok{(fits) <-}\StringTok{ }\KeywordTok{c}\NormalTok{(}\StringTok{"Species"}\NormalTok{,}\StringTok{"c"}\NormalTok{,}\StringTok{"B"}\NormalTok{,}\StringTok{'\textbackslash{}u03B4'}\NormalTok{)}

\KeywordTok{kable}\NormalTok{(fits,}\DataTypeTok{caption=}\StringTok{"Inferred fitness parameters"}\NormalTok{)}
\end{Highlighting}
\end{Shaded}

\begin{table}[t]

\caption{\label{tab:unnamed-chunk-2}Inferred fitness parameters}
\centering
\begin{tabular}{l|r|r|r}
\hline
Species & c & B & δ\\
\hline
M. longirostris & 1.000131 & 0.0522555 & 0.4314043\\
\hline
L. anceps & 44.456503 & 0.1002334 & 0.4314043\\
\hline
\end{tabular}
\end{table}

To visualize how well this inference reflects the data, we plot our
fitted surfaces onto the data in Figure AX. This demonstrates that this
data contains a great deal of noise, but the overall trend is well
captured.

\begin{Shaded}
\begin{Highlighting}[]
\NormalTok{diff <-}\StringTok{ }\KeywordTok{seq}\NormalTok{(}\KeywordTok{min}\NormalTok{(D),}\KeywordTok{max}\NormalTok{(D),}\DataTypeTok{by=}\NormalTok{.}\DecValTok{1}\NormalTok{)}
\KeywordTok{plot}\NormalTok{(D,W1,}\DataTypeTok{main=}\StringTok{"Goodness of fit"}\NormalTok{)}
\KeywordTok{lines}\NormalTok{(diff,}\KeywordTok{W}\NormalTok{(}\OperatorTok{-}\NormalTok{diff,c1,B1,d),}\DataTypeTok{col=}\StringTok{'red'}\NormalTok{)}
\end{Highlighting}
\end{Shaded}

\includegraphics{direct_measurement_files/figure-latex/unnamed-chunk-3-1.pdf}

\begin{Shaded}
\begin{Highlighting}[]
\KeywordTok{plot}\NormalTok{(D,W2)}
\KeywordTok{lines}\NormalTok{(diff,}\KeywordTok{W}\NormalTok{(diff,c2,B2,d),}\DataTypeTok{col=}\StringTok{'blue'}\NormalTok{)}
\end{Highlighting}
\end{Shaded}

\includegraphics{direct_measurement_files/figure-latex/unnamed-chunk-3-2.pdf}

The \emph{R} code used to produce these results is reproduced here:

\begin{Shaded}
\begin{Highlighting}[]
\KeywordTok{require}\NormalTok{(knitr)}
\KeywordTok{require}\NormalTok{(kableExtra)}
\KeywordTok{require}\NormalTok{(xlsx)}
\NormalTok{mamre <-}\StringTok{ }\KeywordTok{read.xlsx}\NormalTok{(}\StringTok{"Mamre interaction for Bob.xlsx"}\NormalTok{,}\DecValTok{1}\NormalTok{)}
\NormalTok{mamre <-}\StringTok{ }\KeywordTok{data.frame}\NormalTok{(mamre}\OperatorTok{$}\NormalTok{nectar_volume_consumed,mamre}\OperatorTok{$}\NormalTok{POLLEN_GRAINS_DEPOSITED,mamre}\OperatorTok{$}\NormalTok{TONGUE.TUBE)}
\KeywordTok{colnames}\NormalTok{(mamre) <-}\StringTok{ }\KeywordTok{c}\NormalTok{(}\StringTok{"nectar consumed"}\NormalTok{,}\StringTok{"pollen grains deposited"}\NormalTok{,}\StringTok{"tongue length - tube length"}\NormalTok{)}
\KeywordTok{kable}\NormalTok{(mamre,}\DataTypeTok{caption=}\StringTok{"Measurements of fitness and morphological differences"}\NormalTok{)}

\NormalTok{W <-}\StringTok{ }\ControlFlowTok{function}\NormalTok{(D,c,B,d)\{}
\NormalTok{  f <-}\StringTok{ }\NormalTok{c}\OperatorTok{*}\KeywordTok{exp}\NormalTok{(}\OperatorTok{-}\NormalTok{B}\OperatorTok{*}\NormalTok{(d}\OperatorTok{+}\NormalTok{D)}\OperatorTok{/}\DecValTok{2}\NormalTok{)}
  \KeywordTok{return}\NormalTok{(f)}
\NormalTok{\}}

\CommentTok{# collect data}
\NormalTok{W1 <-}\StringTok{ }\NormalTok{mamre}\OperatorTok{$}\StringTok{`}\DataTypeTok{nectar consumed}\StringTok{`}
\NormalTok{W2 <-}\StringTok{ }\NormalTok{mamre}\OperatorTok{$}\StringTok{`}\DataTypeTok{pollen grains deposited}\StringTok{`}
\NormalTok{D  <-}\StringTok{ }\NormalTok{mamre}\OperatorTok{$}\StringTok{`}\DataTypeTok{tongue length - tube length}\StringTok{`}
\NormalTok{n <-}\StringTok{ }\KeywordTok{dim}\NormalTok{(mamre)[}\DecValTok{1}\NormalTok{]}

\NormalTok{L <-}\StringTok{ }\ControlFlowTok{function}\NormalTok{(X)\{}
  
  \CommentTok{# unpack parameters}
\NormalTok{  c1 <-}\StringTok{ }\KeywordTok{exp}\NormalTok{(X[}\DecValTok{1}\NormalTok{])}
\NormalTok{  c2 <-}\StringTok{ }\KeywordTok{exp}\NormalTok{(X[}\DecValTok{2}\NormalTok{])}
\NormalTok{  B1 <-}\StringTok{ }\KeywordTok{exp}\NormalTok{(X[}\DecValTok{3}\NormalTok{])}
\NormalTok{  B2 <-}\StringTok{ }\KeywordTok{exp}\NormalTok{(X[}\DecValTok{4}\NormalTok{])}
\NormalTok{  d  <-}\StringTok{ }\KeywordTok{exp}\NormalTok{(X[}\DecValTok{5}\NormalTok{])}
  
\NormalTok{  l <-}\StringTok{ }\DecValTok{0}
  \CommentTok{# accumulate errors across M. longir}
  \ControlFlowTok{for}\NormalTok{(i }\ControlFlowTok{in} \DecValTok{1}\OperatorTok{:}\NormalTok{n)\{}
\NormalTok{    l <-}\StringTok{ }\NormalTok{l }\OperatorTok{+}\StringTok{ }\NormalTok{( }\KeywordTok{W}\NormalTok{(}\OperatorTok{-}\NormalTok{D[i],c1,B1,d) }\OperatorTok{-}\StringTok{ }\NormalTok{W1[i] )}\OperatorTok{^}\DecValTok{2}
\NormalTok{  \}}
  
  \CommentTok{# accumulate errors across L. anceps}
  \ControlFlowTok{for}\NormalTok{(i }\ControlFlowTok{in} \DecValTok{1}\OperatorTok{:}\NormalTok{n)\{}
\NormalTok{    l <-}\StringTok{ }\NormalTok{l }\OperatorTok{+}\StringTok{ }\NormalTok{( }\KeywordTok{W}\NormalTok{(D[i],c2,B2,d) }\OperatorTok{-}\StringTok{ }\NormalTok{W2[i] )}\OperatorTok{^}\DecValTok{2}
\NormalTok{  \}}
  
  \KeywordTok{return}\NormalTok{(}\KeywordTok{sqrt}\NormalTok{(l))}
  
\NormalTok{\}}

\NormalTok{mins <-}\StringTok{ }\OtherTok{NULL}
\NormalTok{vals <-}\StringTok{ }\OtherTok{NULL}
\ControlFlowTok{for}\NormalTok{(i }\ControlFlowTok{in} \DecValTok{1}\OperatorTok{:}\DecValTok{1000}\NormalTok{)\{}
\NormalTok{  out <-}\StringTok{ }\KeywordTok{optim}\NormalTok{(}\KeywordTok{rnorm}\NormalTok{(}\DecValTok{5}\NormalTok{),L)}
\NormalTok{  mins <-}\StringTok{ }\KeywordTok{rbind}\NormalTok{(mins,out}\OperatorTok{$}\NormalTok{par)}
\NormalTok{  vals <-}\StringTok{ }\KeywordTok{c}\NormalTok{(vals,out}\OperatorTok{$}\NormalTok{value)}
\NormalTok{\}}
\NormalTok{i <-}\StringTok{ }\KeywordTok{which.min}\NormalTok{(vals)}
\NormalTok{argmin <-}\StringTok{ }\KeywordTok{exp}\NormalTok{(mins[i,])}

\NormalTok{c1 <-}\StringTok{ }\NormalTok{argmin[}\DecValTok{1}\NormalTok{]}
\NormalTok{c2 <-}\StringTok{ }\NormalTok{argmin[}\DecValTok{2}\NormalTok{]}
\NormalTok{B1 <-}\StringTok{ }\NormalTok{argmin[}\DecValTok{3}\NormalTok{]}
\NormalTok{B2 <-}\StringTok{ }\NormalTok{argmin[}\DecValTok{4}\NormalTok{]}
\NormalTok{d  <-}\StringTok{ }\NormalTok{argmin[}\DecValTok{5}\NormalTok{]}

\NormalTok{fits <-}\StringTok{ }\KeywordTok{data.frame}\NormalTok{(}\KeywordTok{c}\NormalTok{(}\StringTok{"M. longirostris"}\NormalTok{,}\StringTok{"L. anceps"}\NormalTok{),}\KeywordTok{c}\NormalTok{(c1,c2),}\KeywordTok{c}\NormalTok{(B1,B2),}\KeywordTok{c}\NormalTok{(d,d))}
\KeywordTok{colnames}\NormalTok{(fits) <-}\StringTok{ }\KeywordTok{c}\NormalTok{(}\StringTok{"Species"}\NormalTok{,}\StringTok{"c"}\NormalTok{,}\StringTok{"B"}\NormalTok{,}\StringTok{'\textbackslash{}u03B4'}\NormalTok{)}

\KeywordTok{kable}\NormalTok{(fits,}\DataTypeTok{caption=}\StringTok{"Inferred fitness parameters"}\NormalTok{)}

\NormalTok{diff <-}\StringTok{ }\KeywordTok{seq}\NormalTok{(}\KeywordTok{min}\NormalTok{(D),}\KeywordTok{max}\NormalTok{(D),}\DataTypeTok{by=}\NormalTok{.}\DecValTok{1}\NormalTok{)}
\KeywordTok{plot}\NormalTok{(D,W1,}\DataTypeTok{main=}\StringTok{"Goodness of fit"}\NormalTok{)}
\KeywordTok{lines}\NormalTok{(diff,}\KeywordTok{W}\NormalTok{(}\OperatorTok{-}\NormalTok{diff,c1,B1,d),}\DataTypeTok{col=}\StringTok{'red'}\NormalTok{)}
\KeywordTok{plot}\NormalTok{(D,W2)}
\KeywordTok{lines}\NormalTok{(diff,}\KeywordTok{W}\NormalTok{(diff,c2,B2,d),}\DataTypeTok{col=}\StringTok{'blue'}\NormalTok{)}
\end{Highlighting}
\end{Shaded}


\end{document}
